\documentclass[12pt]{article}

% ===================== PACKAGES =====================
\usepackage[margin=1in]{geometry}
\usepackage{graphicx}
\usepackage{listings}
\usepackage{xcolor}
\usepackage{titlesec}
\usepackage{fancyhdr}
\usepackage{hyperref}
\usepackage{array}
\usepackage{booktabs}
\usepackage{tabularx}
\usepackage{float}
\usepackage{enumitem}

% ===================== PAGE STYLE =====================
\pagestyle{fancy}
\fancyhf{}
\rfoot{Page \thepage\ of \pageref{LastPage}}
\usepackage{lastpage}

% ===================== CODE LISTING STYLE =====================
\definecolor{codebg}{rgb}{0.95,0.95,0.95}
\definecolor{codegreen}{rgb}{0,0.5,0}
\definecolor{codegray}{rgb}{0.5,0.5,0.5}
\definecolor{codepurple}{rgb}{0.58,0,0.82}

\lstdefinestyle{cppstyle}{
    backgroundcolor=\color{codebg},
    commentstyle=\color{codegreen},
    keywordstyle=\color{blue}\bfseries,
    stringstyle=\color{codepurple},
    numberstyle=\tiny\color{codegray},
    basicstyle=\ttfamily\footnotesize,
    breakatwhitespace=false,
    breaklines=true,
    captionpos=b,
    keepspaces=true,
    numbers=none,
    showspaces=false,
    showstringspaces=false,
    showtabs=false,
    tabsize=2,
    frame=single,
    language=C++
}
\lstset{style=cppstyle}

\titleformat{\section}{\normalfont\Large\bfseries}{Task \thesection:}{0.5em}{}

% ===================== DOCUMENT =====================
\begin{document}

% ---------- TITLE PAGE ----------
\begin{titlepage}
    \centering
    \vspace*{0.5cm}
    {\itshape Heaven's Light Is Our Guide \par}
    \vspace{0.3cm}
    {\Large \bfseries Rajshahi University of Engineering \& Technology \par}
    \vspace{1cm}

    \includegraphics[width=4cm]{LOGO.png}
    \vspace{1cm}

    {\Large \underline{\bfseries Lab Report 01} \par}
    \vspace{1cm}

    {\large \textbf{Course Title:}  Digital Techniques Sessional \par}
    \vspace{0.3cm}
    {\large \textbf{Course No:} ECE 2112 \par}
    \vspace{0.3cm}
    {\large \textbf{Report Title:} Digital Logic Gates \& Circuits \par}

    \vspace{2cm}

    \noindent
    \begin{tabularx}{\textwidth}{X X}
        \toprule
        \textbf{Submitted By} & \textbf{Submitted To} \\
        \midrule
        Md. Sadikin Naser Araf & Mst. Mazeda Noor Tasnim \\
        Roll: 2410007 & Lecturer \\
        Department of Electrical \& & Department of Electrical \& Computer \\
        Computer Engineering, 24 Series & Engineering, RUET \\
        \bottomrule
    \end{tabularx}

\end{titlepage}



\section*{Introduction}
This report presents a collection of fundamental digital logic circuits implemented and simulated to demonstrate core concepts of digital electronics, including universal gate derivations (NAND/NOR), a Full Adder, and an 8-Bit Binary to BCD converter. All circuits were designed and simulated using Logisim.

\section{All Fundamental Logic Gates using NOR Gate}

The NOR gate is a universal gate, meaning any other logic gate can be constructed using only NOR gates. This circuit demonstrates three separate derivations --- NOT, OR, and AND --- built entirely from NOR gates.

\subsection*{Description}
\begin{itemize}[leftmargin=*]
    \item \textbf{NOT Gate:} Formed by connecting both inputs of a single NOR gate together, so it acts as an inverter.
    \item \textbf{OR Gate:} Formed by feeding two NOR gates in cascade --- the first NOR gate combines A and B, and the second NOR gate (used as a NOT) inverts that result back to a plain OR.
    \item \textbf{AND Gate:} Formed using De Morgan's construction --- each input is first inverted by its own NOR-based NOT gate, and the inverted signals are then combined by a final NOR gate to produce the AND output.
\end{itemize}

\subsubsection*{Truth Table 1 --- NOT Gate (from NOR)}
\begin{table}[H]
\centering
\begin{tabular}{cc}
\toprule
A & Output (NOT A) \\
\midrule
0 & 1 \\
1 & 0 \\
\bottomrule
\end{tabular}
\end{table}

\subsubsection*{Truth Table 2 --- OR Gate (from NOR)}
\begin{table}[H]
\centering
\begin{tabular}{ccc}
\toprule
A & B & Output (A OR B) \\
\midrule
0 & 0 & 0 \\
0 & 1 & 1 \\
1 & 0 & 1 \\
1 & 1 & 1 \\
\bottomrule
\end{tabular}
\end{table}

\subsubsection*{Truth Table 3 --- AND Gate (from NOR)}
\begin{table}[H]
\centering
\begin{tabular}{ccc}
\toprule
A & B & Output (A AND B) \\
\midrule
0 & 0 & 0 \\
0 & 1 & 0 \\
1 & 0 & 0 \\
1 & 1 & 1 \\
\bottomrule
\end{tabular}
\end{table}

\subsection*{Circuit Diagram}
\begin{figure}[H]
\centering
\includegraphics[width=0.7\textwidth]{NORGate.png}

\caption{All Fundamental Logic Gates using NOR Gate}
\end{figure}

\newpage

\section{All Fundamental Logic Gates using NAND Gate}

The NAND gate is also a universal gate. Similar to the NOR-based approach, this circuit demonstrates three separate derivations --- AND, OR, and NOT --- built entirely from NAND gates.

\subsection*{Description}
\begin{itemize}[leftmargin=*]
    \item \textbf{AND Gate:} Formed by cascading two NAND gates --- the first NAND combines A and B, and the second NAND (used as a NOT) inverts that result back to a plain AND.
    \item \textbf{OR Gate:} Formed using De Morgan's construction --- each input is first inverted by its own NAND-based NOT gate, and the inverted signals are then combined by a final NAND gate to produce the OR output.
    \item \textbf{NOT Gate:} Formed by connecting both inputs of a single NAND gate together, so it acts as an inverter.
\end{itemize}

\subsubsection*{Truth Table 1 --- AND Gate (from NAND)}
\begin{table}[H]
\centering
\begin{tabular}{ccc}
\toprule
A & B & Output (A AND B) \\
\midrule
0 & 0 & 0 \\
0 & 1 & 0 \\
1 & 0 & 0 \\
1 & 1 & 1 \\
\bottomrule
\end{tabular}
\end{table}

\subsubsection*{Truth Table 2 --- OR Gate (from NAND)}
\begin{table}[H]
\centering
\begin{tabular}{ccc}
\toprule
A & B & Output (A OR B) \\
\midrule
0 & 0 & 0 \\
0 & 1 & 1 \\
1 & 0 & 1 \\
1 & 1 & 1 \\
\bottomrule
\end{tabular}
\end{table}

\subsubsection*{Truth Table 3 --- NOT Gate (from NAND)}
\begin{table}[H]
\centering
\begin{tabular}{cc}
\toprule
A & Output (NOT A) \\
\midrule
0 & 1 \\
1 & 0 \\
\bottomrule
\end{tabular}
\end{table}

\subsection*{Circuit Diagram}
\begin{figure}[H]
\centering
\includegraphics[width=0.7\textwidth]{Nand GAte.png}

\caption{All Fundamental Logic Gates using NAND Gate}
\end{figure}

\newpage

\section{Full Adder}

A Full Adder is a combinational circuit that performs the addition of three binary bits: two significant bits (A and B) and a carry-in bit (Cin) from a previous addition stage. It produces two outputs: the Sum and the Carry-out (Cout). The circuit is implemented in two ways: using a built-in Adder component (with Cin/Cout), and using two cascaded Half Adders built from XOR/AND/OR gates.

\subsection*{Description}
\begin{itemize}[leftmargin=*]
    \item \textbf{Built-in Full Adder block:} Takes A, B, and Cin as inputs, and directly outputs Sum and Cout.
    \item \textbf{Half Adder 1:} Takes inputs A and B, and produces an intermediate Sum (A XOR B) and an intermediate Carry (A AND B).
    \item \textbf{Half Adder 2 + OR stage:} Takes the intermediate Sum from Half Adder 1 along with Cin to produce the final Sum, while the two intermediate carries are combined with an OR gate to produce the final Cout.
\end{itemize}

\subsubsection*{Truth Table 1 --- Half Adder 1 (A, B)}
\begin{table}[H]
\centering
\begin{tabular}{cccc}
\toprule
A & B & Sum1 (A XOR B) & Carry1 (A AND B) \\
\midrule
0 & 0 & 0 & 0 \\
0 & 1 & 1 & 0 \\
1 & 0 & 1 & 0 \\
1 & 1 & 0 & 1 \\
\bottomrule
\end{tabular}
\end{table}

\subsubsection*{Truth Table 2 --- Half Adder 2 (Sum1, Cin)}
\begin{table}[H]
\centering
\begin{tabular}{cccc}
\toprule
Sum1 & Cin & Sum (Sum1 XOR Cin) & Carry2 (Sum1 AND Cin) \\
\midrule
0 & 0 & 0 & 0 \\
0 & 1 & 1 & 0 \\
1 & 0 & 1 & 0 \\
1 & 1 & 0 & 1 \\
\bottomrule
\end{tabular}
\end{table}

\subsubsection*{Truth Table 3 --- Final Full Adder (A, B, Cin)}
\begin{table}[H]
\centering
\begin{tabular}{ccccc}
\toprule
A & B & Cin & Sum & Cout (Carry1 OR Carry2) \\
\midrule
0 & 0 & 0 & 0 & 0 \\
0 & 0 & 1 & 1 & 0 \\
0 & 1 & 0 & 1 & 0 \\
0 & 1 & 1 & 0 & 1 \\
1 & 0 & 0 & 1 & 0 \\
1 & 0 & 1 & 0 & 1 \\
1 & 1 & 0 & 0 & 1 \\
1 & 1 & 1 & 1 & 1 \\
\bottomrule
\end{tabular}
\end{table}

\subsection*{Circuit Diagram}
\begin{figure}[H]
\centering
\includegraphics[width=0.7\textwidth]{Full Adder.png}

\caption{Full Adder}
\end{figure}

\newpage

\section{8-Bit Binary to BCD Converter}

This circuit converts an 8-bit binary input into its equivalent BCD (Binary Coded Decimal) representation and displays it on three 7-segment displays (Hundreds, Tens, and Units place).

\subsection*{Description}
\begin{itemize}[leftmargin=*]
    \item Takes an 8-bit binary number as input (range: 0--255).
    \item Converts the binary value into three BCD digits: Hundreds (100s), Tens (10s), and Units (1s).
    \item Each BCD digit drives its own 7-segment display, showing the decimal equivalent of the binary input.
    \item Commonly implemented using the Double Dabble (Shift-and-Add-3) algorithm to convert binary to BCD.
    \item Useful in digital systems where binary data needs to be shown in human-readable decimal form, such as counters, digital clocks, and measurement displays.
\end{itemize}

\subsubsection*{Truth Table (Sample Binary to BCD Conversions)}
\begin{table}[H]
\centering
\begin{tabular}{ccccc}
\toprule
Binary Input & Decimal & Hundreds & Tens & Units \\
\midrule
00000000 & 0   & 0 & 0 & 0 \\
00001001 & 9   & 0 & 0 & 9 \\
00011001 & 25  & 0 & 2 & 5 \\
01100100 & 100 & 1 & 0 & 0 \\
11111111 & 255 & 2 & 5 & 5 \\
\bottomrule
\end{tabular}
\end{table}

\subsection*{Circuit Diagram}
\begin{figure}[H]
\centering
\includegraphics[width=0.7\textwidth]{bin to BCD.png}

\caption{8-Bit Binary to BCD Converter}
\end{figure}

\section*{Tools Used}
\begin{itemize}[leftmargin=*]
    \item Logic circuit simulation software (Logisim)
\end{itemize}

\end{document}